% sec_tests_C.tex
% Раздел C (Верификация и тесты): норматив, шлюзы, режимы допуска и витрина протокола
% Актуализировано под рабочий протокол C с QC1 и QC_SCREW_LOCK
% pdfLaTeX-safe: используем LaTeX-математику, без Unicode-математики в обычном тексте
%
% Требует в main.tex:
% \usepackage{amsmath,amssymb}
% \usepackage{tikz}
% \usetikzlibrary{arrows.meta,positioning,shapes.geometric}
% \usepackage{fancyvrb}
% \usepackage{hyperref}

\section{Верификация и тесты (C): норматив, шлюзы и условия допуска к интерпретации}
\label{sec:sec_tests_C}

\subsection{Почему потребовалось ужесточение протокола}

Верификационный контур C был ужесточён по двум независимым причинам.

Во-первых, локально-квазицилиндрическая запись удобна как технический инструмент расчёта, но сама по себе опасна: если использовать её без обязательной парной проверки с глобальной тороидальной постановкой, то можно перепутать \emph{координатный артефакт} с физическим эффектом. Поэтому локальный патч не может служить основанием для вывода без отдельной проверки геометрической инвариантности.

Во-вторых, дискретные шнековые режимы больше нельзя интерпретировать в старом непрерывном $q$-языке. В рабочем интерфейсе B$\to$C физический приоритет перенесён с непрерывного параметра $q$ на дискретную пару
\[
(m_A,m_S)
\]
и метку режима \texttt{mode\_S}. Параметр \texttt{q\_eff} оставлен только как вспомогательный локальный параметр наклона патча и не может служить самостоятельным основанием для физического вывода.

Именно поэтому в протокол C введены два обязательных шлюза:
\begin{itemize}
\item \textbf{QC1} --- защита от координатного артефакта в режиме \texttt{quasi\_cyl};
\item \textbf{QC\_SCREW\_LOCK} --- защита от произвольной непрерывной геометризации шнека.
\end{itemize}

Логика секции C строится в следующем порядке:
\begin{enumerate}
\item сначала фиксируются канонические запреты;
\item затем объясняется, почему \texttt{quasi\_cyl} опасен без пары \texttt{torus\_ref/quasi\_cyl\_ref};
\item затем объясняется, почему шнек нельзя интерпретировать через один \texttt{q\_eff};
\item далее вводятся QC1 и QC\_SCREW\_LOCK как охрана физического смысла;
\item только после этого перечисляются режимы, допускаемые к интерпретации.
\end{enumerate}

\subsection{Канонические запреты и язык наблюдения}

\textbf{(CANON \S16, \S17, норматив)} Прокси-набор $(K,G,A_{ij})$ используется только как
\emph{измерительный язык} состояния хаотической фазы у $\Sigma$ и \emph{не вводит новых сущностей}.
В тексте допустимы формулировки \emph{«наблюдаем прокси»} и запрещены формулировки вида
\emph{«действует поле $G/A$»}, \emph{«тензор $A$ толкает»}, \emph{«$A$ создаёт силу»}.

Рекомендуемая нормировка:
\[
K=n\langle v^2\rangle,\qquad
G=n\langle v\rangle,\qquad
P_{ij}=n\langle v_i v_j\rangle,\qquad
A_{ij}=P_{ij}-\frac13\mathrm{tr}(P)\delta_{ij}.
\]
На поверхности $\Sigma$ используем нормаль $n(x)$ и проектор на касательную плоскость
\[
\Pi_\top(\xi)=\xi-(\xi\cdot n)n.
\]

КУБК относится к балансу \emph{числа контактов} на $\Sigma$ и не является источником дрейфа.
Если приращения импульса изотропны, то устойчивый дрейф запрещён:
\[
A\equiv 0,\qquad \Pi_\top(G)\equiv 0 \quad \Longrightarrow \quad V=0.
\]
Здесь $V=0$ означает \textbf{векторный ноль}, а не только отсутствие продольной компоненты.

\subsection{Охранный тест T2 и нулевой коридор}

\textbf{(CANON \S4.3, \S7, \S17, норматив)} Тест T2 запрещает получать дрейф
из одной лишь $\lambda$-асимметрии. Если заявлено ``нет внешнего нарушителя'' и
входные прокси-метрики находятся в нулевом коридоре, то дрейф обязан оставаться нулевым.

Вводим метрики:
\[
v_{\mathrm{rms}}=\sqrt{\langle K/n\rangle_\Sigma},\qquad \bar n=\langle n\rangle_\Sigma,
\]
\[
\varepsilon_G=\frac{\|\Pi_\top(G)\|_\Sigma}{\bar n\,v_{\mathrm{rms}}},\qquad
\varepsilon_A=\frac{\|A\|_\Sigma}{\bar n\,v_{\mathrm{rms}}^2},
\]
\[
\varepsilon_V=\frac{\|V\|}{v_{\mathrm{rms}}},\qquad
\varepsilon_{V\parallel}=\frac{|V_\parallel|}{v_{\mathrm{rms}}}.
\]

\textbf{Авто-шлюз T2:}
если заявлено отсутствие внешнего нарушителя и
\[
(\varepsilon_G,\varepsilon_A)\approx 0,
\]
то обязано быть
\[
\varepsilon_V\approx 0.
\]

Красный флаг возникает, если $V\neq 0$ при $A\equiv 0$ и $\Pi_\top(G)\equiv 0$:
в этом случае дрейф фактически ``получен из $\lambda$'' или допущена ошибка постановки.

\subsection{Норматив тестов шнека T5--T6}

\textbf{(CANON \S17, норматив)} Фиксируется ось симметрии объекта $e_\parallel$,
измеряется продольная компонента дрейфа:
\[
V_\parallel = V\cdot e_\parallel .
\]

Обязательные проверки:
\[
V_\parallel(\chi)=-V_\parallel(-\chi),
\qquad
V_\parallel(s_{sh}=+1)=-V_\parallel(s_{sh}=-1)
\]
при прочих равных.

Для стоячего шнека при отсутствии внешнего нарушителя:
\[
\nabla K=0,\quad G=0,\quad A=0,\quad s_{sh}=0
\quad \Longrightarrow \quad
V_\parallel\approx 0.
\]

Метрики парных прогонов:
\[
\Delta_\chi=\frac{|V_\parallel(+1)+V_\parallel(-1)|}{v_{\mathrm{rms}}},
\qquad
\Delta_s=\frac{|V_\parallel(s_{sh}=+1)+V_\parallel(s_{sh}=-1)|}{v_{\mathrm{rms}}}.
\]

\subsection{QC1: защита от координатного артефакта \texttt{quasi\_cyl}}

Локальная запись \texttt{quasi\_cyl} допустима только как \emph{временный локальный патч}
той же канонической геометрии. Она не является новой физической сущностью и не заменяет
глобальную тороидальную/двуториальную постановку.

Поэтому для каждой канонической ситуации B обязан формировать минимум одну парную связку:
\begin{itemize}
\item \texttt{geometry\_mode=global\_torus}, \texttt{pair\_role=torus\_ref};
\item \texttt{geometry\_mode=quasi\_cyl}, \texttt{pair\_role=quasi\_cyl\_ref}.
\end{itemize}

При этом обязаны совпадать:
\begin{itemize}
\item \texttt{scenario},
\item \texttt{geometry\_id},
\item \texttt{pair\_id\_geom},
\item физические входы,
\item \texttt{chi},
\item \texttt{s\_sh},
\item режим внешнего нарушителя,
\item настройки допусков,
\end{itemize}
и различаться только:
\begin{itemize}
\item \texttt{geometry\_mode},
\item \texttt{pair\_role}.
\end{itemize}

\textbf{QC1 PASS:}
\begin{itemize}
\item охранные запреты T2, N1, N2, N3 не ломаются при переходе к \texttt{quasi\_cyl};
\item новый дрейф не возникает только из выбора координатного патча;
\item допустимость/недопустимость D1 совпадает в паре \texttt{torus\_ref/quasi\_cyl\_ref}.
\end{itemize}

\textbf{QC1 FAIL:}
\begin{itemize}
\item запреты нарушаются только в \texttt{quasi\_cyl};
\item новый эффект есть только в \texttt{quasi\_cyl}, а в \texttt{global\_torus} исчезает.
\end{itemize}

В этом случае серия маркируется как
\[
\texttt{coordinate\_artifact\_suspected},
\]
а физическая интерпретация откладывается.

\subsection{QC\_SCREW\_LOCK: защита от произвольной непрерывной геометризации шнека}

В новой рабочей схеме B$\to$C физика шнека читается через дискретную пару
\[
(m_A,m_S),
\]
где:
\begin{itemize}
\item $m_A\in \mathbb{Z}_{>0}$ --- квантовое число канальной навивки;
\item $m_S\in\{1,2,4\}$ --- базовая шнековая мода;
\item \texttt{mode\_S} --- метка режима:
\[
\texttt{mode\_S}\in\{\texttt{S1},\texttt{S2},\texttt{S4},\texttt{S*}\}.
\]
\end{itemize}

Базовые режимы:
\[
\texttt{S1},\ \texttt{S2},\ \texttt{S4}.
\]
Режим \texttt{S*} допускается только как \emph{производный} режим с явной маркировкой производности.

Дополнительно B обязан передавать:
\begin{itemize}
\item \texttt{hA\_definition},
\item \texttt{hS\_definition},
\item \texttt{A\_S\_orth\_candidate},
\item \texttt{q\_eff},
\item \texttt{mode\_lock\_pass}.
\end{itemize}

\textbf{Критический смысл:} \texttt{q\_eff} не является главным физическим носителем шнека.
Он допустим только как вспомогательный локальный параметр патча.

\textbf{QC\_SCREW\_LOCK PASS:}
\begin{itemize}
\item серия читается через согласованную пару $(m_A,m_S)$;
\item базовые режимы принадлежат множеству \texttt{S1/S2/S4};
\item \texttt{S*} указан как производный режим;
\item \texttt{h\_A} и \texttt{h\_S} разведены явно;
\item \texttt{q\_eff} не используется как физическая основа режима;
\item \texttt{mode\_lock\_pass=true}.
\end{itemize}

\textbf{QC\_SCREW\_LOCK FAIL:}
\begin{itemize}
\item шнековая мода подана как произвольная непрерывная деформация;
\item связь с $m_A$ отсутствует;
\item режим вне базы \texttt{S1/S2/S4} заявлен без указания производности;
\item физическая интерпретация даётся при \texttt{mode\_lock\_pass=false};
\item \texttt{q\_eff} используется как главный аргумент физического вывода.
\end{itemize}

\subsection{Базовые моды S1/S2/S4 и условия допуска к интерпретации}

Для каждого нового шнекового сценария C требует:
\begin{enumerate}
\item парность \texttt{global\_torus} $\leftrightarrow$ \texttt{quasi\_cyl};
\item сравнение базовых мод \texttt{S1}, \texttt{S2}, \texttt{S4};
\item сравнение ``базовая мода $\leftrightarrow$ производная мода'';
\item сравнение ``согласованная пара $(m_A,m_S)$ $\leftrightarrow$ искусственно рассогласованная пара''.
\end{enumerate}

Физическая интерпретация шнековой серии допустима только если одновременно:
\begin{itemize}
\item \texttt{local\_patch\_valid=true};
\item \texttt{patch\_valid=true} (если используется патч-язык);
\item \texttt{QC1\_pass=true};
\item \texttt{QC\_SCREW\_LOCK\_pass=true};
\item \texttt{mode\_lock\_pass=true};
\item пройден T2;
\item выполнены T5--T6;
\item пройдены N1--N3;
\item D1 заявляется только в допустимом режиме.
\end{itemize}

Если хотя бы одно из условий нарушено, C обязан выставить:
\[
\texttt{FAIL\_INTERPRETATION}
\]
и/или соответствующий красный флаг.

\subsection{Охранники N1--N3 и режим D1 в новой схеме}

Перед тем, как считать режим D1 физически интерпретируемым, обязательны охранники:

\begin{itemize}
\item \textbf{N1 (базовая тороидальная ссылка):}
нет допустимого дискретно-согласованного пространственного режима; локальный tor-вклад и вращение допустимы,
но интегрально обязано быть
\[
V\approx 0.
\]

\item \textbf{N2 (без tor-моды):}
\[
s_{sh}=0 \quad \Longrightarrow \quad
j_{\mathrm{tor}}^{\mathrm{loc}}=0
\Longrightarrow
j_{\mathrm{sp}}^{\mathrm{loc}}=0
\Longrightarrow
V=0.
\]

\item \textbf{N3 (без геометрического рычага):}
\[
w_{\mathrm{asym}}\approx 0 \quad \Longrightarrow \quad V\approx 0.
\]
\end{itemize}

\textbf{D1 (дискретно-согласованный двуторий со шнеком)} допускается только если:
\begin{itemize}
\item предъявлена допустимая пара $(m_A,m_S)$;
\item \texttt{mode\_S}$\in\{\texttt{S1},\texttt{S2},\texttt{S4}\}$ или явно обоснованная производная \texttt{S*};
\item $w_{\mathrm{asym}}>0$;
\item $s_{sh}\neq 0$.
\end{itemize}

Только здесь допускается
\[
V_\parallel\neq 0,
\]
но с обязательным выполнением норм CANON \S17:
\[
V_\parallel(\chi)=-V_\parallel(-\chi),
\qquad
V_\parallel(s_{sh}=+1)=-V_\parallel(s_{sh}=-1).
\]

\subsection{Почему в витрине должны быть две поточные схемы}

В этой секции намеренно сохраняются \textbf{две} визуальные схемы.

Первая схема --- \emph{историческая}. Она нужна не как действующий норматив, а как
документированная фиксация прошлого этапа, в котором C строил объяснение через
непрерывную геометризацию двутория при помощи $q$, $h(x)$ и $w(x)$.
Эта схема показывает, откуда произошёл ранний язык охранников N1--N3 и D1.

Вторая схема --- \emph{актуальная}. Она нужна как действующий поток проверки в текущем
рабочем протоколе C, где центральны уже не параметры $q/h(x)$, а:
\begin{itemize}
\item \texttt{geometry\_mode},
\item временный локальный патч,
\item QC1,
\item дискретная пара $(m_A,m_S)$,
\item \texttt{mode\_S},
\item QC\_SCREW\_LOCK,
\item и только затем --- физическая интерпретация.
\end{itemize}

Сохранение обеих схем нужно для честной демонстрации всех шагов эволюции протокола:
мы не скрываем прежний каркас, но и не выдаём его за действующий.

\subsection{Историческая поточная схема C (ранняя q-центричная витрина)}

\begin{figure}[p]
\centering
\resizebox{\textwidth}{0.90\textheight}{%
\begin{tikzpicture}[
  font=\normalsize,
  node distance=9mm and 24mm,
  box/.style={rectangle, rounded corners=3pt, draw=black!55, thick, align=center,
              minimum width=0.78\linewidth, inner sep=8pt},
  sbox/.style={rectangle, rounded corners=3pt, draw=black!55, thick, align=center,
               minimum width=0.36\linewidth, inner sep=8pt},
  decision/.style={diamond, aspect=2.25, draw=black!55, thick, align=center,
                   inner sep=3pt, text width=0.44\linewidth},
  arr/.style={-Latex, thick, draw=black!70}
]

\node[box, fill=blue!10] (geom) {%
$\Sigma$: граничная поверхность, параметризация $X(u,v)$\\
$n(x)$: нормаль на $\Sigma$,\quad $e_{\parallel}$: ось (CANON \S17),\quad $V_{\parallel}=V\cdot e_{\parallel}$\\
паспорт: $\chi\in\{-1,+1\}$,\quad $s_{sh}\in\{-1,0,+1\}$ (бег $\rightarrow$/стоячий/бег $\leftarrow$)
};

\node[box, fill=cyan!10, below=of geom] (geo) {%
\textbf{Геометризация двутория (B$\to$C, КАНДИДАТ):}\\
$q\in\mathbb{R}\setminus\{0\}$ (q=0: тонкий тор; q$\neq$0: двуторий),\ \ $w(x)$: вес внешн./внутр. рычага\\
\texttt{w\_mode $\in$ \{binary, hybrid, ...\}},\quad \texttt{w\_asym} — метрика несимметрии $w(x)$\\
Норма каркаса: дрейф двутория запрещён без (q$\neq$0 и w\_asym$>$0) и без tor-моды (s\_sh$\neq$0)
};

\node[box, fill=orange!12, below=of geo] (proxies) {%
Прокси на $\Sigma$ (CANON \S16):\\
$K=n\langle v^2\rangle,\quad G=n\langle v\rangle,\quad P_{ij}=n\langle v_i v_j\rangle$\\
$A_{ij}=P_{ij}-\frac13\mathrm{tr}(P)\delta_{ij},\quad \Pi_{\top}(G)=G-(G\!\cdot\!n)n$\\
\textbf{Прокси — ярлыки измерения, не сущности и не «поля».}
};

\node[decision, fill=white, below=of proxies] (dec) {Расчёт $j$ на $\Sigma$};

\node[sbox, fill=violet!12, below left=of dec, yshift=+4mm] (ext) {%
\textbf{Внешний}\\[2pt]
$j^{ext}(x)=F_{ext}(K,G,A;\Sigma)$\\
$\nabla K$ / приходящая ВР\\
(перекосы статистики на $\Sigma$)
};

\node[sbox, fill=violet!12, below right=of dec, yshift=+4mm] (intr) {%
\textbf{Внутренний (tor$\to$sp)}\\[2pt]
$g_{\top}=\Pi_{\top}(G),\quad a_{\top}=\Pi_{\top}(A n)$\\
$j_{tor}=\chi\,s_{sh}(\alpha_1\,n\times g_{\top}+\alpha_2\,n\times a_{\top})$\\
$j_{sp}=\beta(K)\,w(x)\,(j_{tor}\!\cdot\!t_\theta)\,h(x)$\\
$j_{int}=j_{tor}+j_{sp}$\\
\textbf{s\_sh=0 $\Rightarrow$ $j_{tor}=j_{sp}=0$}
};

\node[box, fill=green!12, below=30mm of dec] (sum) {%
$j(x)=j^{ext}(x)+j^{int}(x)$\\
локальная импульсная анизотропия контактов на $\Sigma$
};

\node[box, fill=purple!10, below=of sum] (resp) {%
Интегральный отклик:\\
$V=\mu_T\int_{\Sigma} j\,dS,\quad V_{\parallel}=V\cdot e_{\parallel}$\\
$\Omega=\tilde{\mu}_R\int_{\Sigma}(r-r_0)\times j\,dS$\\
$J_{int,\parallel}=\Big(\int_{\Sigma} j^{int}\,dS\Big)\cdot e_{\parallel}$
};

\node[box, fill=yellow!18, below=of resp] (ver) {%
Верификация / нормы:\\
$v_{rms}=\sqrt{\langle K/n\rangle_{\Sigma}},\quad \bar n=\langle n\rangle_{\Sigma}$\\
$\varepsilon_G=\|\Pi_{\top}(G)\|_{\Sigma}/(\bar n\,v_{rms}),\quad
\varepsilon_A=\|A\|_{\Sigma}/(\bar n\,v_{rms}^2)$\\
$\varepsilon_V=\|V\|/v_{rms},\quad \varepsilon_{V\parallel}=|V_{\parallel}|/v_{rms}$\\
\textbf{T2:} $(\varepsilon_G,\varepsilon_A)\approx 0 \Rightarrow \varepsilon_V\approx 0$\\
\textbf{T5--T6:} $\chi$-нечётность $V_{\parallel}$ и реверс $s_{sh}$; $s_{sh}=0 \Rightarrow V_{\parallel}\approx 0$\\[4pt]
\textbf{Охранники геометризации (перед D1 обязательно):}\\
N1: $q=0,\ s_{sh}\neq 0 \Rightarrow \Omega\neq 0$ допустимо, но $V\approx 0$\\
N2: $q\neq 0,\ s_{sh}=0 \Rightarrow V=0$\\
N3: $q\neq 0,\ w\_{asym}\approx 0 \Rightarrow V\approx 0$\\
\textbf{D1:} $q\neq 0,\ w\_{asym}>0,\ s_{sh}\neq 0 \Rightarrow V_{\parallel}\neq 0$ + нормы CANON \S17
};

\draw[arr] (geom) -- (geo);
\draw[arr] (geo) -- (proxies);
\draw[arr] (proxies) -- (dec);
\draw[arr] (dec) -- (ext);
\draw[arr] (dec) -- (intr);
\draw[arr] (ext.south)  -- ++(0,-10mm) |- (sum.west);
\draw[arr] (intr.south) -- ++(0,-10mm) |- (sum.east);
\draw[arr] (sum) -- (resp);
\draw[arr] (resp) -- (ver);

\end{tikzpicture}%
}
\caption{Историческая поточная схема (C): ранняя витрина через непрерывную геометризацию двутория $(q,w,h)$ и переход $tor\to sp$. Схема сохраняется как документированная версия прежнего этапа развития протокола.}
\label{fig:C_flow_tikz_old}
\end{figure}

\subsection{Актуальная поточная схема C (действующий протокол)}

Новая схема вводится не для косметической замены обозначений, а для изменения самой логики проверки.
Если историческая схема показывала, как C раньше описывал внутренний механизм через непрерывную геометризацию,
то актуальная схема показывает, как C теперь проверяет \emph{допустимость физического вывода}.

В новой логике:
\begin{enumerate}
\item локальный патч признаётся только временным инструментом расчёта;
\item затем обязательна геометрическая проверка QC1;
\item затем обязательна модовая проверка QC\_SCREW\_LOCK;
\item только после этого допускается интерпретация результата как физического.
\end{enumerate}

Именно поэтому в новой схеме:
\begin{itemize}
\item появился шаг \emph{временного локального патча};
\item появился явный шлюз \textbf{QC1};
\item появился явный шлюз \textbf{QC\_SCREW\_LOCK};
\item параметр $q$ ушёл с роли главного физического носителя шнека;
\item интерпретация вынесена в последний блок потока.
\end{itemize}

\begin{figure}[p]
\centering
% Новый блок на ТОЧНО ТОМ ЖЕ геометрическом каркасе, что и old:
% меняется только содержание узлов, а не внешняя компоновка.
\resizebox{\textwidth}{0.90\textheight}{%
\begin{tikzpicture}[
  font=\normalsize,
  node distance=9mm and 24mm,
  box/.style={rectangle, rounded corners=3pt, draw=black!55, thick, align=center,
              minimum width=0.78\linewidth, inner sep=8pt},
  sbox/.style={rectangle, rounded corners=3pt, draw=black!55, thick, align=center,
               minimum width=0.36\linewidth, inner sep=8pt},
  decision/.style={diamond, aspect=2.25, draw=black!55, thick, align=center,
                   inner sep=3pt, text width=0.44\linewidth},
  arr/.style={-Latex, thick, draw=black!70}
]

% --- 1) Геометрия
\node[box, fill=blue!10] (geom) {%
$\Sigma$: граничная поверхность, параметризация $X(u,v)$\\
$n(x)$: нормаль на $\Sigma$,\quad $e_{\parallel}$: ось (CANON \S17),\quad $V_{\parallel}=V\cdot e_{\parallel}$\\
паспорт: $\chi\in\{-1,+1\}$,\quad $s_{sh}\in\{-1,0,+1\}$,\quad \texttt{geometry\_mode $\in$ \{global\_torus, quasi\_cyl\}}
};

% --- 2) Временный локальный патч
\node[box, fill=cyan!10, below=of geom] (geo) {%
\textbf{Временный локальный патч и геометрический режим}\\
\texttt{global\_torus}: глобальная тороидальная / двуториальная постановка\\
\texttt{quasi\_cyl}: только локальный патч записи $\Sigma_{loc}\subset\Sigma$\\
\texttt{local\_patch\_valid}, \texttt{patch\_valid}, \texttt{curvature\_scale\_ratio}, \texttt{curvature\_ratio}\\
\textbf{Норма:} \texttt{quasi\_cyl} не является новой физической сущностью
};

% --- 3) Прокси
\node[box, fill=orange!12, below=of geo] (proxies) {%
Прокси на $\Sigma$ (CANON \S16):\\
$K=n\langle v^2\rangle,\quad G=n\langle v\rangle,\quad P_{ij}=n\langle v_i v_j\rangle$\\
$A_{ij}=P_{ij}-\frac13\mathrm{tr}(P)\delta_{ij},\quad \Pi_{\top}(G)=G-(G\!\cdot\!n)n$\\
\textbf{Прокси --- ярлыки измерения, не сущности и не «поля».}
};

% --- 4) Ромб: QC1
\node[decision, fill=white, below=of proxies] (dec) {%
\textbf{QC1}\\
геометрическая инвариантность\\
\texttt{torus\_ref} $\leftrightarrow$ \texttt{quasi\_cyl\_ref}
};

% --- 5) Внешний
\node[sbox, fill=violet!12, below left=of dec, yshift=+4mm] (ext) {%
\textbf{Внешний}\\[2pt]
$j^{ext}(x)=F_{ext}(K,G,A;\Sigma)$\\
$\nabla K$ / приходящая ВР\\
внешний перекос статистики на $\Sigma$
};

% --- 6) Внутренний: новая модовая схема
\node[sbox, fill=violet!12, below right=of dec, yshift=+4mm] (intr) {%
\textbf{Внутренний (дискретно-модовый)}\\[2pt]
$m_A\in\mathbb{Z}_{>0},\quad m_S\in\{1,2,4\}$,\quad \texttt{mode\_S $\in$ \{S1,S2,S4,S*\}}\\
$h_A=h_A(m_A;x),\quad h_S=h_S(m_S,\chi,s_{sh};x)$\\
$j_{tor}^{loc}=\chi\,s_{sh}\,\mathcal{T}(m_A;K,g_{\top},a_{\top},n)$\\
$j_{sp}^{loc}=\beta(K;m_A,m_S)\,w_{loc}(x)\,\Gamma_{A,S}(x)\,h_S(x)$\\
\textbf{\texttt{q\_eff} --- только вспомогательный параметр патча}
};

% --- 7) Сумма / модовый шлюз
\node[box, fill=green!12, below=30mm of dec] (sum) {%
\textbf{QC\_SCREW\_LOCK и суммарный отклик}\\
база: \texttt{S1/S2/S4},\quad \texttt{S*} только как производная,\quad \texttt{mode\_lock\_pass=true}\\
$j(x)=j^{ext}(x)+j^{int}(x)$\\
локальная импульсная анизотропия контактов на $\Sigma$
};

% --- 8) Интегральный отклик
\node[box, fill=purple!10, below=of sum] (resp) {%
Интегральный отклик:\\
$V=\mu_T\int_{\Sigma} j\,dS,\quad V_{\parallel}=V\cdot e_{\parallel}$\\
$\Omega=\tilde{\mu}_R\int_{\Sigma}(r-r_0)\times j\,dS$\\
$J_{int,\parallel}=\Big(\int_{\Sigma} j^{int}\,dS\Big)\cdot e_{\parallel}$
};

% --- 9) Нормы/шлюзы + N1–N3 + D1
\node[box, fill=yellow!18, below=of resp] (ver) {%
Верификация / нормы:\\
$v_{rms}=\sqrt{\langle K/n\rangle_{\Sigma}},\quad \bar n=\langle n\rangle_{\Sigma}$\\
$\varepsilon_G=\|\Pi_{\top}(G)\|_{\Sigma}/(\bar n\,v_{rms}),\quad
\varepsilon_A=\|A\|_{\Sigma}/(\bar n\,v_{rms}^2)$\\
$\varepsilon_V=\|V\|/v_{rms},\quad \varepsilon_{V\parallel}=|V_{\parallel}|/v_{rms}$\\
\textbf{T2:} $(\varepsilon_G,\varepsilon_A)\approx 0 \Rightarrow \varepsilon_V\approx 0$\\
\textbf{T5--T6:} $\chi$-нечётность $V_{\parallel}$ и реверс $s_{sh}$; $s_{sh}=0 \Rightarrow V_{\parallel}\approx 0$\\[4pt]
\textbf{Охранники перед D1:}\\
N1: базовая тороидальная ссылка $\Rightarrow V\approx 0$\\
N2: $s_{sh}=0 \Rightarrow j_{tor}^{loc}=0 \Rightarrow j_{sp}^{loc}=0 \Rightarrow V=0$\\
N3: $w\_{asym}\approx 0 \Rightarrow V\approx 0$\\
\textbf{D1:} допустим только при корректной паре $(m_A,m_S)$ и режиме \texttt{S1/S2/S4} или производной \texttt{S*}
};

% --- Стрелки
\draw[arr] (geom) -- (geo);
\draw[arr] (geo) -- (proxies);
\draw[arr] (proxies) -- (dec);
\draw[arr] (dec) -- (ext);
\draw[arr] (dec) -- (intr);

% Входы в sum по серединам граней (как в old):
\draw[arr] (ext.south)  -- ++(0,-10mm) |- (sum.west);
\draw[arr] (intr.south) -- ++(0,-10mm) |- (sum.east);

\draw[arr] (sum) -- (resp);
\draw[arr] (resp) -- (ver);

\end{tikzpicture}%
}
\caption{Актуальная поточная схема (C): сохранён прежний внешний каркас витрины, но внутренняя логика обновлена под режимы \texttt{global\_torus/quasi\_cyl}, геометрический шлюз QC1, дискретно-модовую схему шнека и модовый шлюз QC\_SCREW\_LOCK.}
\label{fig:C_flow_tikz_new}
\end{figure}

\FloatBarrier

\subsection{Реестр красных флагов}

\begin{itemize}
\item дрейф появляется при изотропном фоне без нарушителя;
\item дрейф ``объясняется'' асимметрией числа контактов, а не анизотропией $\Delta p$;
\item $\Phi$ интерпретируется как поток СЭ;
\item $(K,G,A_{ij})$ трактуются как действующие поля/сущности;
\item \texttt{quasi\_cyl} трактуется как отдельная физическая сущность или глобальная форма ЭВО;
\item новый эффект существует только в \texttt{quasi\_cyl};
\item \texttt{q\_eff} используется как главный физический носитель шнека;
\item режим вне базы \texttt{S1/S2/S4} интерпретируется без явной производности;
\item физическая интерпретация дана при \texttt{mode\_lock\_pass=false};
\item D1 заявлен без прохождения N1--N3, QC1 и QC\_SCREW\_LOCK.
\end{itemize}
